\documentclass[10pt,a4paper]{amsart}
\usepackage[margin=19mm]{geometry}
\usepackage{amsmath,amssymb,mathtools}
\usepackage[scr=rsfs,cal=euler]{mathalfa}
\usepackage{xfrac}
\usepackage[unicode,hidelinks,bookmarks=false]{hyperref}
\newcommand{\ie}{i.e.\ }
\newcommand{\cf}{cf.\ }
\newcommand{\resp}{respectively\ }
\newcommand{\Subsection}{\subsection}
\providecommand{\nobreakdash}{\nobreak\hbox{-}}
\setlength{\emergencystretch}{2em}
\multlinegap0pt
\usepackage[shortlabels]{enumitem}
\usepackage{amssymb}
\usepackage{mathtools}
\usepackage{float}
\newtheorem{corollary}{Corollary}
\newtheorem{theorem}{Theorem}
\usepackage{cleveref}
\crefname{corollary}{Corollary}{Corollarys}
\crefname{theorem}{Theorem}{Theorems}
\title{Inverse Theorems of Point--Sphere Incidences over Finite Fields}
\author{Shalender Singh, Vishnu Priya Singh Parmar}
\date{Source: arXiv:2602.10123v3 (CC BY 4.0)}
\begin{document}
\maketitle

\noindent\emph{Source and licence.} Inverse Theorems of Point--Sphere Incidences over Finite Fields. Version arXiv:2602.10123v3 (CC BY 4.0), distributed by arXiv under the Creative Commons Attribution 4.0 International licence, http://creativecommons.org/licenses/by/4.0/, as recorded in the arXiv licence metadata for this exact version and shown on the abstract page https://arxiv.org/abs/2602.10123v3. Full licence notice: \texttt{licenses/arxiv-2602.10123-CC-BY-4.0.txt}.\par\smallskip

\noindent\textit{Editorial scope.} Counted unit: the article's theorem thm:rnc-block-indecomp (source0.tex lines 1821--1841). The article's complete original proof of it is delivered with it (source0.tex lines 1843--1935, one \texttt{proof} environment of 93 lines). No mathematics is written, repaired or re-derived: the statement and the proof are the article's own lines, in the article's own notation, and no retained line is substituted or re-flowed.  Retained uncounted as the source-local context the proof needs: lines 872--887; lines 903--916; lines 1618--1622; lines 1624--1673; lines 1633--1638; lines 1811--1817.  Nothing was omitted from any retained span, so the package carries no omission record.  No retained line contains a citation, so the wrapper carries no bibliography and no bibliography entry.  No retained line contains a figure environment, an image-inclusion command or drawing code, so no image is delivered and the package declares no asset dependency.  Editorial formatting only, in addition: an AMS \texttt{amsart} preamble that restates the article's own declarations and macros used by the retained lines, namely the environments \texttt{corollary}, \texttt{theorem} on separate counters, together with the standard packages those lines need.  The retained environments print in this document as Corollary 1, Theorem 1 in order of appearance, whereas the article prints the same items with the numbering of its own full document.  The complete native source of the article as distributed in the arXiv e-print for this exact version is bundled unchanged as \texttt{sources/194371/source0.tex}.
\begin{corollary}[Fixed-scale adaptive inverse theorem]\label{cor:fixed-adaptive}
Let $P\subseteq\mathbb F_q^d$ be nonempty, let $\mathcal S$ be a nonempty family of distinct spheres, and let $K>0$.  Suppose
\begin{equation}\label{eq:fixed-surplus}
 \sigma(P,\mathcal S)\ge Kq^{(d-1)/2}\sqrt{|P||\mathcal S|}
\end{equation}
and
\begin{equation}\label{eq:point-sensitive-lower}
 K^2q^{d-1}|\mathcal S|\ge4|P|.
\end{equation}
Then at least $K^2|\mathcal S|/4$ ordered pairs satisfy
\begin{equation}\label{eq:adaptive-section}
 |P\cap S\cap S'|\ge
 \frac{K^2q^{d-1}}{4|\mathcal S|}.
\end{equation}
Every such pair supplies a radical hyperplane containing at least the same number of points.  At least $K^2/8$ spheres each have at least $K^2/8$ witnessing partners.  If $d\ge3$, at least $K^2q|\mathcal S|/8$ ordered pairs satisfy \eqref{eq:adaptive-section}, and at least $K^2q/16$ spheres each have at least $K^2q/16$ witnessing partners.
\end{corollary}

\begin{corollary}[Moderate-family inverse theorem]\label{cor:moderate}
Let $P\subseteq\mathbb F_q^d$ be nonempty, let $\mathcal S$ be a nonempty family of distinct spheres, let $K>0$, and let $C_1\ge1$.  Suppose \eqref{eq:fixed-surplus} holds and
\begin{equation}\label{eq:moderate-window}
 \frac{4|P|}{K^2q^{d-1}}
 \le|\mathcal S|
 \le C_1Kq^{(d-1)/2}.
\end{equation}
Then at least $K^2|\mathcal S|/4$ ordered pairs of spheres satisfy
\begin{equation}\label{eq:moderate-section}
 |P\cap S\cap S'|\ge
 \frac{Kq^{(d-1)/2}}{4C_1}.
\end{equation}
Their radical hyperplanes contain at least the same number of points.  If $d\ge3$, the number of witnessing pairs is at least $K^2q|\mathcal S|/8$.  The partner-abundance conclusions of \cref{cor:fixed-adaptive} hold as well.
\end{corollary}

\begin{equation}\label{eq:rnc-system}
 P_{\mathrm{rn}}=\{p_A:A\in\mathfrak A_{d,\gamma}\},
 \qquad
 \mathcal S_{\mathrm{rn}}=\{S_t:t\in\mathbb F_q\}.
\end{equation}

\begin{theorem}[Diffuse rational-normal-curve sphere systems]\label{thm:rnc-system}
Fix $d\ge3$.  For all sufficiently large odd $q$, the system in \eqref{eq:rnc-system} has the following properties.  Write $N=|P_{\mathrm{rn}}|$ and $\Sigma=\sigma(P_{\mathrm{rn}},\mathcal S_{\mathrm{rn}})$.
\begin{enumerate}[label=\textup{(\arabic*)}]
\item Every sphere has nonzero radius.  The lifted parameters $\Phi(S_t)$ are affinely equivalent to
\[
 t\longmapsto(t,t^2,\ldots,t^d)\in\mathbb F_q^d.
\]
For every $0\le k\le d-1$, no affine $k$-flat in their $d$-dimensional affine span contains more than $k+1$ lifted parameters.  In particular, no coaxal pencil contains three members of $\mathcal S_{\mathrm{rn}}$.
\item The map $A\mapsto p_A$ is injective,
\begin{equation}\label{eq:rnc-N-bounds}
 \binom qd-\frac1d\binom q{d-1}
 \le N\le\binom qd,
 \qquad
 N=\frac{q^d}{d!}+O_d(q^{d-1}),
\end{equation}
and
\begin{equation}\label{eq:rnc-incidence}
 p_A\in S_t\quad\Longleftrightarrow\quad t\in A.
\end{equation}
Consequently every point has degree $d$.  For distinct $t,u$,
\begin{equation}\label{eq:rnc-pair-count}
 \binom{q-2}{d-2}-\frac1{d-2}\binom{q-2}{d-3}
 \le |P_{\mathrm{rn}}\cap S_t\cap S_u|
 \le\binom{q-2}{d-2}.
\end{equation}
\item One has
\begin{equation}\label{eq:rnc-surplus}
 I(P_{\mathrm{rn}},\mathcal S_{\mathrm{rn}})=dN,
 \qquad
 \Sigma=(d-1)N,
 \qquad
 K_q=\frac{(d-1)\sqrt N}{q^{d/2}}
 \longrightarrow\frac{d-1}{\sqrt{d!}}.
\end{equation}
Moreover,
\begin{equation}\label{eq:rnc-lower-exact}
 K_q^2q^{d-1}|\mathcal S_{\mathrm{rn}}|=(d-1)^2N\ge4N.
\end{equation}
For all sufficiently large $q$, the system lies in the moderate window of \cref{cor:moderate} with $C_1=2$, and its rich-pair graph is complete.
\item Every $B\subseteq P_{\mathrm{rn}}$ satisfies
\begin{equation}\label{eq:rnc-uniform-excess}
 \sigma(B,\mathcal S_{\mathrm{rn}})=(d-1)|B|.
\end{equation}
Hence, uniformly over distinct $t,u$,
\begin{equation}\label{eq:rnc-no-atom}
 \frac{\sigma(P_{\mathrm{rn}}\cap S_t\cap S_u,\mathcal S_{\mathrm{rn}})}{\Sigma}
 =O_d(q^{-2}).
\end{equation}
\end{enumerate}
\end{theorem}

\begin{equation}
 \binom qd-\frac1d\binom q{d-1}
 \le N\le\binom qd,
 \qquad
 N=\frac{q^d}{d!}+O_d(q^{d-1}),
\end{equation}

\begin{align}
 \sigma(A_i,\mathcal S_i)&\ge\tau\Sigma
 &&(1\le i\le m),\label{eq:block-diagonal}\\
 \sigma(R,\mathcal S)&\le\varepsilon\Sigma,\label{eq:block-remainder}\\
 \sum_{i\ne j}\sigma(A_i,\mathcal S_j)_+&\le\varepsilon\Sigma,
 \qquad x_+=\max\{x,0\}.\label{eq:block-cross}
\end{align}

\begin{theorem}[Quantitative centered-block indecomposability]\label{thm:rnc-block-indecomp}
Fix $d\ge3$ and put
\[
 c_d=\frac{d-1}{32d^2},
 \qquad
 C_d=8d(d+1).
\]
Uniformly for every $\gamma\in\mathbb F_q$, every $0<\tau\le1$, and every odd $q\ge C_d/\tau$, the root-evaluation system $\mathcal D_{d,\gamma}$ admits no $(\tau,\varepsilon)$-centered block decomposition with
\[
 \varepsilon\le c_d\tau^2.
\]
In particular, this conclusion holds for the rational-normal-curve sphere system of \cref{thm:rnc-system}.  More quantitatively, writing $P$ and $\mathcal S$ for the two sides of $\mathcal D_{d,\gamma}$ and $\Sigma=\sigma(P,\mathcal S)$, every matched partition satisfying \eqref{eq:block-diagonal} and
\[
 \sigma(R,\mathcal S)\le c_d\tau^2\Sigma
\]
obeys
\begin{equation}\label{eq:rnc-cross-lower}
 \sum_{i\ne j}\sigma(A_i,\mathcal S_j)_+
 \ge2c_d\tau^2\Sigma.
\end{equation}
\end{theorem}

\begin{proof}
Identify each object with its label $t\in\mathbb F_q$ and each point with its $d$-element incidence neighborhood.  Put $N=|P|$.  The deletion argument used in the proof of \eqref{eq:rnc-N-bounds} is uniform in $\gamma$, and gives
\[
 \binom qd-\frac1d\binom q{d-1}\le N\le\binom qd.
\]
Every point has degree $d$, $|\mathcal S|=q$, and therefore $\Sigma=(d-1)N$.  Write
\[
 s_i=|\mathcal S_i|,
 \qquad
 \theta_i=\frac{s_i}{q},
 \qquad
 n_i=|A_i|,
\]
and let $n_i^0$ be the number of points of $A_i$ whose entire incidence neighborhood lies in $\mathcal S_i$.  Put
\[
 \mu=\frac{d!N}{q^d},
 \qquad
 \rho=\frac{\tau(d-1)}{2d}.
\]
The lower bound for $N$ gives the explicit estimate
\begin{align}
 \mu
 &\ge \prod_{j=0}^{d-1}\left(1-\frac jq\right)-\frac1q\notag\\
 &\ge1-\frac{d(d-1)/2+1}{q}
 \ge1-\frac\rho4,
 \label{eq:rnc-mu}
\end{align}
where the last inequality follows from $q\ge C_d/\tau$ and $d(d-1)/2+1\le(d+1)(d-1)$.

A fixed object belongs to at most $\binom{q-1}{d-1}\le q^{d-1}/(d-1)!$ points.  If $D_i=\sigma(A_i,\mathcal S_i)$, then
\[
 D_i\le I(P,\mathcal S_i)
 \le\frac{d\theta_iq^d}{d!}.
\]
On the other hand, \eqref{eq:block-diagonal} and $\Sigma=(d-1)N$ give
\[
 D_i\ge\tau(d-1)N
 =\frac{\tau(d-1)\mu q^d}{d!}.
\]
Thus $\theta_i\ge2\rho\mu\ge\rho$.  Since $m\ge2$, also $\theta_i\le1-\rho$ for every $i$.

For a point $e\in A_i$, let $k_i(e)$ be the number of its $d$ incident object labels lying in $\mathcal S_i$.  Summing the signed cross discrepancies leaving the point block $A_i$ gives
\[
 C_i:=\sum_{j\ne i}\sigma(A_i,\mathcal S_j)
 =\sum_{e\in A_i}\bigl((d-k_i(e))-(1-\theta_i)\bigr).
\]
A point internal to $\mathcal S_i$ contributes $-(1-\theta_i)$, whereas every other point contributes at least $\theta_i$.  Therefore
\begin{equation}\label{eq:rnc-row-cross}
 C_i\ge\theta_i n_i-n_i^0.
\end{equation}
Put
\[
 X=\sum_{i\ne j}\sigma(A_i,\mathcal S_j)_+.
\]
Using \eqref{eq:rnc-row-cross}, $\theta_i\ge\rho$, and the elementary inequality $\sum_j(x_j)_+\ge(\sum_jx_j)_+$, we obtain
\begin{align}
 X
 &\ge\sum_i(C_i)_+
 \ge\sum_i(\theta_i n_i-n_i^0)_+\notag\\
 &\ge\rho\sum_i\left(n_i-\frac{n_i^0}{\theta_i}\right)_+
 \ge\rho\left(\sum_i n_i-\sum_i\frac{n_i^0}{\theta_i}\right)_+.
 \label{eq:rnc-X}
\end{align}

Now
\[
 n_i^0\le\binom{s_i}{d}\le\frac{\theta_i^dq^d}{d!},
\]
so
\begin{align*}
 \sum_i\frac{n_i^0}{\theta_i}
 &\le\frac{q^d}{d!}\sum_i\theta_i^{d-1}
 \le\frac{q^d}{d!}(1-\rho)^{d-2}\\
 &\le\frac{q^d}{d!}(1-\rho)
 \le\left(1-\frac\rho2\right)N.
\end{align*}
The last inequality follows from $\mu\ge1-\rho/4$ and
\[
 \frac{1-\rho}{1-\rho/4}\le1-\frac\rho2.
\]
Since every point has full-family excess degree $d-1$, the remainder hypothesis implies
\[
 |R|\le c_d\tau^2N\le\frac\rho4N;
\]
the last inequality holds because $0<\tau\le1$ and $d\ge3$.
Thus $\sum_i n_i=N-|R|\ge(1-\rho/4)N$.  Substitution in \eqref{eq:rnc-X} yields
\[
 X\ge\frac{\rho^2}{4}N
 =\frac{\tau^2(d-1)}{16d^2}\Sigma
 =2c_d\tau^2\Sigma,
\]
which proves \eqref{eq:rnc-cross-lower} and the theorem.
\end{proof}

\end{document}
